\paragraph{软核张量算子的全谱分解、割线方程与惯性平衡}
设 \(q\ge2\)，\(U_q:=\{0,1\}^q\)，并按标准张量识别
\[
\mathbb R^{U_q}\cong(\mathbb R^2)^{\otimes q}.
\]
记
\[
K:=\begin{pmatrix}1&1\\[2pt]1&0\end{pmatrix},
\qquad
J:=\begin{pmatrix}1&1\\[2pt]1&1\end{pmatrix},
\]
\[
D^{(q)}:=K^{\otimes q},
\qquad
J_q:=J^{\otimes q},
\qquad
T^{(q)}:=\frac12\bigl(J_q+D^{(q)}\bigr).
\]
令
\[
\varphi:=\frac{1+\sqrt5}{2},
\qquad
\psi:=\frac{1-\sqrt5}{2}=-\varphi^{-1},
\]
\[
\mu_i^{(q)}:=\varphi^{\,q-i}\psi^{\,i}=(-1)^i\varphi^{\,q-2i},
\qquad 0\le i\le q.
\]
并记 \(\mathcal H_q:=\mathrm{Sym}^q(\mathbb R^2)\subset\mathbb R^{U_q}\) 为 \(S_q\)-对称子空间。

\begin{theorem}[全谱的例外--体谱直和分解与精确重数]\label{thm:xi-replica-softcore-full-spectrum-direct-sum}
对每个 \(0\le i\le q\)，\(\mu_i^{(q)}/2\) 是 \(T^{(q)}\) 的特征值，其代数重数恰为
\[
\binom{q}{i}-1.
\]
其余恰好 \(q+1\) 个特征值为 \(\left.T^{(q)}\right|_{\mathcal H_q}\) 的谱
\(\{\lambda_0^{(q)},\dots,\lambda_q^{(q)}\}\)，并与体谱不相交。

等价地，存在正交直和分解
\[
\mathbb R^{U_q}
=
\mathcal H_q\ \oplus\
\bigoplus_{i=0}^{q}\bigl(E_i\cap\mathbf 1^\perp\bigr),
\]
其中 \(E_i:=\ker(D^{(q)}-\mu_i^{(q)}I)\)、\(\mathbf1\) 为全 \(1\) 向量，且
\[
\left.T^{(q)}\right|_{E_i\cap\mathbf1^\perp}
=
\frac{\mu_i^{(q)}}{2}I.
\]
\end{theorem}
\begin{proof}
矩阵 \(K\) 为实对称矩阵，特征值为 \(\varphi,\psi\)，故 \(D^{(q)}=K^{\otimes q}\) 的特征值为
\(\mu_i^{(q)}\)，且
\[
\dim E_i=\binom{q}{i}.
\]
又 \(J_q\) 为秩一算子且
\[
\operatorname{Ran}(J_q)=\operatorname{span}\{\mathbf1\},
\qquad
J_qx=0\iff x\perp\mathbf1.
\]
因此对任意 \(x\in E_i\cap\mathbf1^\perp\)，
\[
T^{(q)}x=\frac12\bigl(D^{(q)}x+J_qx\bigr)=\frac{\mu_i^{(q)}}2x.
\]

记 \(x_i:=\operatorname{Proj}_{E_i}\mathbf1\)。
由于 \((1,1)\) 在 \(K\) 的两条特征方向上投影均非零，故 \(x_i\neq0\) 对全部 \(i\) 成立，从而
\[
\dim(E_i\cap\mathbf1^\perp)=\dim E_i-1=\binom{q}{i}-1.
\]
体谱总维数为
\[
\sum_{i=0}^{q}\left(\binom{q}{i}-1\right)=2^q-(q+1).
\]
余下 \(q+1\) 维空间由 \(\{x_i\}_{i=0}^{q}\) 张成；该空间 \(S_q\)-不变且维数与 \(\mathcal H_q\) 相同，故即为 \(\mathcal H_q\)。
由结论 \ref{con:xi-replica-softcore-exceptional-secular-strict-interlacing}，\(\mathcal H_q\) 上的例外谱与 \(\{\mu_i^{(q)}/2\}\) 严格交错，因而与体谱不相交。证毕。
\end{proof}

\begin{corollary}[特征多项式分解与体谱因子闭式]\label{cor:xi-replica-softcore-charpoly-bulk-factorization}
记
\[
\chi_q(x):=\det(xI-T^{(q)}),
\qquad
\chi_q^{\mathrm{exc}}(x):=\det\!\left(xI-\left.T^{(q)}\right|_{\mathcal H_q}\right).
\]
则
\[
\chi_q(x)
=
\chi_q^{\mathrm{exc}}(x)\,
\prod_{i=0}^{q}\left(x-\frac{\mu_i^{(q)}}2\right)^{\binom{q}{i}-1}.
\]
若记 \(m_q^{\mathrm{exc}}\) 为 \(\left.T^{(q)}\right|_{\mathcal H_q}\) 的最小多项式，则
\[
m_q(x)=m_q^{\mathrm{exc}}(x)\prod_{i=0}^{q}\left(x-\frac{\mu_i^{(q)}}2\right).
\]
\end{corollary}
\begin{proof}
由定理 \ref{thm:xi-replica-softcore-full-spectrum-direct-sum} 的不变直和分解，特征多项式按分块相乘。
又体谱点在对应分块上为纯标量，指数即其重数。
最小多项式由各不变分量最小多项式的最小公倍式给出，并由谱分离得到所示因子。证毕。
\end{proof}

\begin{theorem}[例外谱的加权割线方程]\label{thm:xi-replica-softcore-exceptional-weighted-secular}
定义
\[
\alpha:=1+\frac{2}{\sqrt5},
\qquad
\beta:=1-\frac{2}{\sqrt5},
\]
\[
w_i^{(q)}:=\binom{q}{i}\alpha^{\,q-i}\beta^{\,i}>0,
\qquad 0\le i\le q.
\]
则 \(\left.T^{(q)}\right|_{\mathcal H_q}\) 的全部特征值恰为方程
\[
\sum_{i=0}^{q}\frac{w_i^{(q)}}{\lambda-\mu_i^{(q)}/2}=2
\]
的全部根。
并且
\[
\sum_{i=0}^{q}w_i^{(q)}=2^q,
\]
所有根均为单根。
\end{theorem}
\begin{proof}
取 \(K\) 的单位正交本征向量 \(e_+,e_-\)（对应 \(\varphi,\psi\)），并记
\[
c_+:=\langle(1,1),e_+\rangle,\qquad c_-:=\langle(1,1),e_-\rangle.
\]
直接计算得
\[
c_+^2=1+\frac{2}{\sqrt5}=\alpha,\qquad
c_-^2=1-\frac{2}{\sqrt5}=\beta.
\]
在 \(\mathcal H_q\) 的标准对称本征基 \(\{s_i\}_{i=0}^{q}\) 下，
\[
\left.D^{(q)}\right|_{\mathcal H_q}=\operatorname{diag}\bigl(\mu_0^{(q)},\dots,\mu_q^{(q)}\bigr),
\]
且 \(\mathbf1\) 的坐标向量 \(g\) 满足
\[
g_i^2=\binom{q}{i}c_+^{2(q-i)}c_-^{2i}=w_i^{(q)}.
\]
故
\[
\left.T^{(q)}\right|_{\mathcal H_q}
=
\operatorname{diag}\!\left(\frac{\mu_i^{(q)}}2\right)_{i=0}^{q}
+\frac12\,gg^{\mathsf T}.
\]
由矩阵行列式引理，
\[
0=
\det\!\left(
\lambda I-\operatorname{diag}\!\left(\frac{\mu_i^{(q)}}2\right)-\frac12gg^{\mathsf T}
\right)
\]
\[
=
\prod_{i=0}^{q}\left(\lambda-\frac{\mu_i^{(q)}}2\right)
\left(
1-\frac12\sum_{i=0}^{q}\frac{g_i^2}{\lambda-\mu_i^{(q)}/2}
\right).
\]
又 \(g_i\neq0\)，故极点 \(\lambda=\mu_i^{(q)}/2\) 不属例外谱，从而得到割线方程。
\(\sum_i w_i^{(q)}=(\alpha+\beta)^q=2^q\) 由二项式定理即得。
根的单性由
\(
\frac{\mathrm d}{\mathrm d\lambda}\sum_i\frac{w_i^{(q)}}{\lambda-\mu_i^{(q)}/2}<0
\)
推出。证毕。
\end{proof}

\begin{proposition}[秩一割线方程的耦合单调性与导数--范数同一律]\label{prop:xi-rankone-secular-monotonicity-derivative-norm}
\leanverified{paper\_xi\_rankone\_secular\_monotonicity\_derivative\_norm}
取互异实数 \(t_0,\dots,t_q\) 与权重 \(w_i>0\)，并令
\[
D:=\mathrm{diag}(t_0,\dots,t_q),\qquad g:=(\sqrt{w_0},\dots,\sqrt{w_q})^{\mathsf T}.
\]
对任意 \(b>0\)，考虑秩一扰动
\[
A(b):=D+b\,gg^{\mathsf T}.
\]
则 \(A(b)\) 的谱由割线方程
\begin{equation}\label{eq:xi-rankone-secular-general}
\Phi(\lambda):=\sum_{i=0}^{q}\frac{w_i}{\lambda-t_i}=\frac1b,
\qquad
\lambda\notin\{t_0,\dots,t_q\},
\end{equation}
给出，且任意根分支 \(\lambda=\lambda(b)\) 在其定义域内满足严格单调与导数闭式
\begin{equation}\label{eq:xi-rankone-secular-derivative}
\frac{\mathrm d\lambda}{\mathrm db}
=\frac{1}{b^2\displaystyle\sum_{i=0}^{q}\frac{w_i}{(\lambda-t_i)^2}}
\;>\;0.
\end{equation}
并且与 \(\lambda(b)\) 对应的特征向量可取为
\[
v(b):=(\lambda I-D)^{-1}g,\qquad
v_i(b)=\frac{\sqrt{w_i}}{\lambda(b)-t_i},
\]
从而耦合量与范数满足恒等式
\[
g^{\mathsf T}v(b)=\frac1b,
\qquad
\|v(b)\|_2^2=\sum_{i=0}^{q}\frac{w_i}{(\lambda(b)-t_i)^2}=\frac{1}{b^2\,\lambda'(b)}.
\]
\end{proposition}
\begin{proof}
由矩阵行列式引理，
\[
\det(\lambda I-A(b))
=\det(\lambda I-D)\Bigl(1-b\,g^{\mathsf T}(\lambda I-D)^{-1}g\Bigr),
\]
其中
\(
g^{\mathsf T}(\lambda I-D)^{-1}g=\sum_i\frac{w_i}{\lambda-t_i}=\Phi(\lambda)
\)。
故 \(\lambda\) 为特征值当且仅当 \eqref{eq:xi-rankone-secular-general} 成立；并且 \(\Phi'(\lambda)=-\sum_i\frac{w_i}{(\lambda-t_i)^2}<0\)，根均为单根。
对 \(\Phi(\lambda(b))=1/b\) 作隐式微分即得 \eqref{eq:xi-rankone-secular-derivative}，从而 \(\lambda'(b)>0\)。
\par
取 \(v(b)=(\lambda I-D)^{-1}g\)，则
\[
(A(b)-\lambda I)v
=(D-\lambda I)v+b\,g(g^{\mathsf T}v)
=-g+b\,g\,\Phi(\lambda)
=0,
\]
其中最后一步使用割线方程；故 \(v(b)\) 为特征向量。
其坐标式与 \(g^{\mathsf T}v=1/b\) 由定义直接给出；
范数恒等式由坐标式与 \eqref{eq:xi-rankone-secular-derivative} 联立得到。
\end{proof}

\begin{theorem}[秩一耦合的谱纤维分裂：小耦合与大耦合极限]\label{thm:xi-rankone-secular-small-large-coupling}
沿用命题 \ref{prop:xi-rankone-secular-monotonicity-derivative-norm} 的记号。
令极点集合 \(t_i\) 两两不同，并记
\[
G(\lambda):=\prod_{i=0}^{q}(\lambda-t_i),
\qquad
N(\lambda):=\sum_{i=0}^{q} w_i\!\!\prod_{\substack{0\le j\le q\\ j\ne i}}\!\!(\lambda-t_j),
\qquad
\Phi(\lambda)=\frac{N(\lambda)}{G(\lambda)}.
\]
则对 \(b>0\) 的全部特征值（即 \(\Phi(\lambda)=1/b\) 的全部根）有：
\begin{enumerate}
  \item 当 \(b\to 0^+\) 时，对每个 \(i\in\{0,\dots,q\}\) 存在唯一根分支 \(\lambda_i(b)\) 满足
  \[
  \lambda_i(b)=t_i+bw_i+O(b^2),
  \qquad
  \lambda_i(b)\to t_i.
  \]
  \item 当 \(b\to+\infty\) 时，存在唯一根分支 \(\lambda_\infty(b)\to+\infty\) 且
  \[
  \lambda_\infty(b)=b\sum_{i=0}^{q}w_i+\frac{\sum_{i=0}^{q}w_it_i}{\sum_{i=0}^{q}w_i}+O(b^{-1}).
  \]
  其余 \(q\) 条根分支在 \(b\to+\infty\) 时分别收敛到多项式 \(N(\lambda)\) 的 \(q\) 个实零点；这些零点均为单根并严格夹在相邻极点之间。
\end{enumerate}
\end{theorem}
\begin{proof}
\par\noindent\textbf{(i) 小耦合.}
固定 \(i\)。在 \(\lambda=t_i\) 的穿孔邻域内，
\[
\Phi(\lambda)=\frac{w_i}{\lambda-t_i}+H_i(\lambda),
\qquad
H_i(\lambda):=\sum_{j\ne i}\frac{w_j}{\lambda-t_j},
\]
其中 \(H_i\) 在 \(\lambda=t_i\) 解析。
令 \(\lambda=t_i+\delta\)，方程 \(\Phi(\lambda)=1/b\) 等价于
\(
w_i/\delta+H_i(t_i)+O(\delta)=1/b
\)。
当 \(b\to 0^+\) 时右端趋于 \(+\infty\)，迫使 \(\delta\to 0\)，并由上式解得
\(\delta=bw_i+O(b^2)\)。
由于 \(\Phi'(\lambda)<0\)，在足够小的邻域内根唯一，从而得到唯一根分支。
\par
\noindent\textbf{(ii) 大耦合.}
当 \(|\lambda|\to\infty\)，有渐近展开
\[
\Phi(\lambda)=\sum_{i=0}^{q}\frac{w_i}{\lambda-t_i}
=\frac{W}{\lambda}+\frac{M_1}{\lambda^2}+O(\lambda^{-3}),
\qquad
W:=\sum_{i=0}^{q}w_i,\quad M_1:=\sum_{i=0}^{q}w_it_i.
\]
令 \(\Phi(\lambda)=1/b\) 并令 \(b\to+\infty\) 得到唯一逃逸根满足 \(\lambda\sim bW\)；
进一步迭代代回即得所示渐近展开。
\par
其余根的收敛：注意 \(N(\lambda)=0\) 当且仅当 \(\Phi(\lambda)=0\)。
在每个相邻极点之间的开区间上，\(\Phi\) 连续且严格单调（\(\Phi'<0\)），并在端点趋于 \(\pm\infty\)，因此 \(N\) 在每个此类区间内恰有一个实零点且为单根。
同理，\(\Phi(\lambda)=1/b\) 在每个区间内亦恰有一个实根；令 \(b\to+\infty\) 使右端趋于 \(0\)，由单调性与连续性得该根收敛到区间内的零点。
\end{proof}

\endinput


\begin{theorem}[二项--几何谱测度的 Fibonacci 幂矩塌缩]\label{thm:xi-replica-softcore-fibonacci-moment-collapse}
沿用定理 \ref{thm:xi-replica-softcore-exceptional-weighted-secular} 的权重 \(w_i^{(q)}\) 与节点
\(
d_i:=\mu_i^{(q)}/2=\frac12\varphi^{\,q-i}\psi^{\,i}=\frac12(-1)^i\varphi^{\,q-2i}
\)（\(0\le i\le q\)）。
则对任意整数 \(m\ge 0\)，成立严格恒等式
\begin{equation}\label{eq:xi-replica-softcore-fibonacci-moment-collapse}
\sum_{i=0}^{q} w_i^{(q)}\,d_i^{\,m}
=2^{-m}\,F_{m+3}^{\,q}.
\end{equation}
因此所有幂矩均在 \(\QQ\) 内闭合；右端不含任何 \(\sqrt5\) 的显式项。
\end{theorem}
\begin{proof}
由 \(\psi=-\varphi^{-1}\)，有
\(
d_i^{\,m}=2^{-m}\varphi^{m(q-i)}\psi^{mi}
\)。
因此
\[
\sum_{i=0}^{q} w_i^{(q)}\,d_i^{\,m}
=2^{-m}\sum_{i=0}^{q}\binom{q}{i}\alpha^{q-i}\beta^{i}\varphi^{m(q-i)}\psi^{mi}
=2^{-m}\sum_{i=0}^{q}\binom{q}{i}(\alpha\varphi^m)^{q-i}(\beta\psi^m)^{i}.
\]
二项式定理给出
\[
\sum_{i=0}^{q} w_i^{(q)}\,d_i^{\,m}
=2^{-m}\bigl(\alpha\varphi^m+\beta\psi^m\bigr)^{q}.
\]
再由
\(
\alpha=\varphi^3/\sqrt5
\)、\(
\beta=-\psi^3/\sqrt5
\)，得到
\[
\alpha\varphi^m+\beta\psi^m=\frac{\varphi^{m+3}-\psi^{m+3}}{\sqrt5}=F_{m+3},
\]
代回即得 \eqref{eq:xi-replica-softcore-fibonacci-moment-collapse}。证毕。
\end{proof}
\leanverified{paper\_xi\_replica\_softcore\_fibonacci\_moment\_collapse}

\paragraph{对称幂谱的 Vandermonde 判别式与秩一耦合结果式}
在本段，Fibonacci 幂矩塌缩 \eqref{eq:xi-replica-softcore-fibonacci-moment-collapse} 与秩一割线方程链（命题 \ref{prop:xi-rankone-secular-monotonicity-derivative-norm}）联结为一条纯代数的判别式/结果式闭式，并进一步给出 Hankel 行列式的同型桥接。

\begin{theorem}[Fibonacci--Vandermonde 判别式分解]\label{thm:xi-symq-k-discriminant-fibonacci-vandermonde}
\leanverified{paper\_xi\_symq\_k\_discriminant\_fibonacci\_vandermonde}
令
\[
K=\begin{pmatrix}1&1\\[2pt]1&0\end{pmatrix},
\qquad
B_q:=\mathrm{Sym}^q(K)\in M_{q+1}(\ZZ),
\qquad
g_q(z):=\det(zI_{q+1}-B_q).
\]
则对一切 \(q\ge 2\)，成立严格恒等式
\[
\Disc(g_q)
=5^{\frac{q(q+1)}2}\prod_{k=1}^{q}F_k^{\,2(q+1-k)}.
\]
\end{theorem}
\begin{proof}
由于 \(K\) 为实对称矩阵，其特征值为 \(\varphi\) 与 \(\psi=-\varphi^{-1}\)；于是
\(\mathrm{Sym}^q(K)\) 的谱为
\[
\mu_i=\varphi^{\,q-i}\psi^{\,i}=\varphi^{\,q}\,r^i,
\qquad
r:=\psi/\varphi=-\varphi^{-2},
\qquad
0\le i\le q.
\]
从而 \(g_q(z)=\prod_{i=0}^{q}(z-\mu_i)\) 为首一多项式，并有
\[
\Disc(g_q)=\prod_{0\le i<j\le q}(\mu_j-\mu_i)^2.
\]
注意
\(
\mu_j-\mu_i=\varphi^{\,q}r^i(r^{j-i}-1)
\)，故 Vandermonde 乘积可分解为
\[
\prod_{0\le i<j\le q}(\mu_j-\mu_i)
=\varphi^{\,q\sum_{i=0}^{q}(q-i)}\,
r^{\sum_{i=0}^{q}i(q-i)}\,
\prod_{k=1}^{q}(r^k-1)^{q+1-k}.
\]
平方后得
\[
\Disc(g_q)
=\varphi^{\,2q\sum_{i=0}^{q}(q-i)}\,
r^{2\sum_{i=0}^{q}i(q-i)}\,
\prod_{k=1}^{q}(r^k-1)^{2(q+1-k)}.
\]
其中
\[
\sum_{i=0}^{q}(q-i)=\frac{q(q+1)}2,
\qquad
\sum_{i=0}^{q}i(q-i)=\frac{q(q+1)(q-1)}6.
\]
又因 \(r=-\varphi^{-2}\) 且 \(\psi^k=(-1)^k\varphi^{-k}\)，Binet 恒等式给出
\[
1-(-\varphi^{-2})^k
=\frac{\varphi^k-\psi^k}{\varphi^k}
=\frac{\sqrt5\,F_k}{\varphi^k},
\]
从而
\[
(r^k-1)^2=(1-(-\varphi^{-2})^k)^2=\frac{5F_k^2}{\varphi^{2k}}.
\]
代入并合并指数可见所有 \(\varphi\)-幂严格消去，于是得到所示闭式。证毕。
\end{proof}

\begin{theorem}[秩一耦合的 Sylvester 结果式：谱差积的 Fibonacci 分解]\label{thm:xi-symq-kj-resultant-fibonacci}\leanverified{paper\_xi\_symq\_kj\_resultant\_fibonacci}
令
\[
J=\begin{pmatrix}1&1\\[2pt]1&1\end{pmatrix},
\qquad
E_q:=\mathrm{Sym}^q(J),
\qquad
A_q:=B_q+E_q,
\qquad
p_q(z):=\det(zI_{q+1}-A_q).
\]
则对一切 \(q\ge 2\)，成立严格恒等式
\[
\operatorname{Res}_z(g_q,p_q)
=(-1)^{\frac{(q+1)(q+2)}2}
\left(\prod_{i=0}^{q}\binom{q}{i}\right)
\left(\prod_{k=1}^{q}F_k^{\,2(q+1-k)}\right).
\]
\end{theorem}
\begin{proof}
记 \(s:=(1,1)^{\mathsf T}\)。由于 \(J=ss^{\mathsf T}\)，有
\[
E_q=\mathrm{Sym}^q(J)=u\,u^{\mathsf T},
\qquad
u:=s^{\otimes q}\in\mathrm{Sym}^q(\RR^2).
\]
取 \(B_q\) 的一组标准正交特征基 \(\{e_i\}_{i=0}^{q}\) 使 \(B_q e_i=\mu_i e_i\)。
将 \(u\) 展开为 \(u=\sum_{i=0}^{q}u_i e_i\)，则在该特征基下 \(A_q=\mathrm{diag}(\mu_i)+uu^{\mathsf T}\)。
由矩阵行列式引理得对任意 \(z\) 有
\[
p_q(z)=g_q(z)\Bigl(1-\sum_{i=0}^{q}\frac{u_i^2}{z-\mu_i}\Bigr).
\]
因此对每个根 \(\mu_i\) 有
\[
p_q(\mu_i)=-u_i^2\,g_q'(\mu_i).
\]
由结果式的根值表述，
\[
\operatorname{Res}_z(g_q,p_q)=\prod_{i=0}^{q}p_q(\mu_i)
=(-1)^{q+1}\Bigl(\prod_{i=0}^{q}u_i^2\Bigr)\Bigl(\prod_{i=0}^{q}g_q'(\mu_i)\Bigr).
\]
又因 \(g_q(z)=\prod_{i=0}^{q}(z-\mu_i)\) 首一，
\[
\prod_{i=0}^{q}g_q'(\mu_i)
=\prod_{i=0}^{q}\prod_{j\ne i}(\mu_i-\mu_j)
=(-1)^{\frac{q(q+1)}2}\Disc(g_q).
\]
为计算 \(\prod u_i^2\)，取 \(K\) 的单位正交本征向量 \(e_+,e_-\)（对应 \(\varphi,\psi\)），并记
\[
c_+:=\langle s,e_+\rangle,\qquad c_-:=\langle s,e_-\rangle.
\]
则由对称幂的二项展开，在 \(\mathrm{Sym}^q(\RR^2)\) 的标准正交对称本征基下有
\[
u_i=\sqrt{\binom{q}{i}}\,c_+^{\,q-i}c_-^{\,i},
\qquad 0\le i\le q.
\]
故
\[
\prod_{i=0}^{q}u_i^2
=\left(\prod_{i=0}^{q}\binom{q}{i}\right)
c_+^{\,2\sum_{i=0}^{q}(q-i)}c_-^{\,2\sum_{i=0}^{q}i}
=\left(\prod_{i=0}^{q}\binom{q}{i}\right)(c_+c_-)^{q(q+1)}.
\]
而 \(c_+^2=1+\frac{2}{\sqrt5}\)、\(c_-^2=1-\frac{2}{\sqrt5}\)，从而
\((c_+c_-)^2=\frac15\)，即 \(c_+c_-=5^{-1/2}\)。
因此
\[
\prod_{i=0}^{q}u_i^2
=\left(\prod_{i=0}^{q}\binom{q}{i}\right)5^{-\frac{q(q+1)}2}.
\]
将其与定理 \ref{thm:xi-symq-k-discriminant-fibonacci-vandermonde} 的 \(\Disc(g_q)\) 代回，并合并符号指数
\[
(-1)^{q+1+\frac{q(q+1)}2}=(-1)^{\frac{(q+1)(q+2)}2},
\]
即得所示结果式恒等式。证毕。
\end{proof}

\begin{proposition}[Hankel--Sylvester 桥接：谱差结果式的幂矩行列式模型]\label{prop:xi-symq-hankel-sylvester-bridge}\leanverified{paper\_xi\_symq\_hankel\_sylvester\_bridge}
令
\[
H_q:=\bigl(F_{i+j+3}^{\,q}\bigr)_{0\le i,j\le q}.
\]
则对一切 \(q\ge 2\)，成立
\[
\operatorname{Res}_z(g_q,p_q)
=(-1)^{\frac{(q+1)(q+2)}2}\det(H_q).
\]
\end{proposition}
\begin{proof}
由幂矩塌缩 \eqref{eq:xi-replica-softcore-fibonacci-moment-collapse}，对任意 \(n\ge 0\) 有
\[
F_{n+3}^{\,q}=2^{n}\sum_{i=0}^{q}w_i^{(q)}\,d_i^{\,n},
\qquad
d_i=\mu_i/2,
\]
其中 \(w_i^{(q)}>0\) 与 \(d_i\) 见定理 \ref{thm:xi-replica-softcore-fibonacci-moment-collapse}。
因此
\[
H_q
=D\cdot \Bigl(\sum_{i=0}^{q}w_i^{(q)}\,v(d_i)v(d_i)^{\mathsf T}\Bigr)\cdot D,
\qquad
D:=\mathrm{diag}(2^0,2^1,\dots,2^{q}),
\]
其中 \(v(d):=(1,d,\dots,d^{q})^{\mathsf T}\)。
由 Gram--Vandermonde 行列式公式得到
\[
\det(H_q)
=2^{q(q+1)}\Bigl(\prod_{i=0}^{q}w_i^{(q)}\Bigr)\prod_{0\le i<j\le q}(d_j-d_i)^2.
\]
将 \(d_i=\mu_i/2\) 代入并整理得
\[
\det(H_q)
=\Bigl(\prod_{i=0}^{q}w_i^{(q)}\Bigr)\prod_{0\le i<j\le q}(\mu_j-\mu_i)^2
=\Bigl(\prod_{i=0}^{q}w_i^{(q)}\Bigr)\Disc(g_q).
\]
而 \(\prod_{i=0}^{q}w_i^{(q)}=\bigl(\prod_{i=0}^{q}\binom{q}{i}\bigr)\alpha^{\frac{q(q+1)}2}\beta^{\frac{q(q+1)}2}\)
且 \(\alpha\beta=\frac15\)，从而
\(\prod w_i^{(q)}=\bigl(\prod_{i=0}^{q}\binom{q}{i}\bigr)5^{-\frac{q(q+1)}2}\)。
结合定理 \ref{thm:xi-symq-k-discriminant-fibonacci-vandermonde} 即得
\[
\det(H_q)
=\left(\prod_{i=0}^{q}\binom{q}{i}\right)\left(\prod_{k=1}^{q}F_k^{\,2(q+1-k)}\right).
\]
与定理 \ref{thm:xi-symq-kj-resultant-fibonacci} 对比即得所示桥接恒等式。证毕。
\end{proof}

\begin{theorem}[Hankel 秩与线性递推：最小阶的可检验刻画]\label{thm:xi-hankel-rank-linear-recurrence}
令 \((a_n)_{n\ge 0}\) 为复数列，定义无限 Hankel 矩阵
\[
H=\bigl(a_{i+j}\bigr)_{i,j\ge 0},
\]
以及主子式
\[
\Delta_m:=\det\bigl(a_{i+j}\bigr)_{0\le i,j\le m}\qquad(m\ge 0).
\]
则下列条件等价：
\begin{enumerate}
  \item 存在整数 \(k\ge 1\) 与常数 \(c_0,\dots,c_{k-1}\in\CC\)，使得对一切 \(n\ge 0\) 有
  \[
  a_{n+k}=\sum_{j=0}^{k-1}c_j\,a_{n+j};
  \]
  \item \(\mathrm{rank}(H)\le k\)；
  \item 对所有 \(m\ge k\) 有 \(\Delta_m=0\)。
\end{enumerate}
并且最小可行递推阶数恰等于 \(\mathrm{rank}(H)\)（等价地：恰为使得 \(\Delta_m=0\) 对所有 \(m\ge k\) 成立的最小 \(k\)）。
\leanverified{paper\_xi\_hankel\_rank\_linear\_recurrence}
\end{theorem}
\begin{proof}
\((1)\Rightarrow(2)\)：递推给出对每个 \(j\ge k\)，Hankel 的第 \(j\) 列是前 \(k\) 列的线性组合，故秩至多 \(k\)。
\((2)\Rightarrow(3)\)：秩至多 \(k\) 等价于任意 \((k+1)\times(k+1)\) 子式为零，尤其所有主子式 \(\Delta_m\) 在 \(m\ge k\) 时为零。
\((3)\Rightarrow(1)\)：若所有 \((k+1)\times(k+1)\) 主子式为零，则 Hankel 前 \(k+1\) 列线性相关，存在不全为零的 \((\gamma_0,\dots,\gamma_k)\) 使
\[
\sum_{j=0}^k \gamma_j a_{n+j}=0\qquad(\forall n\ge 0),
\]
从而得到阶至多 \(k\) 的递推。最小阶数等于 Hankel 秩由列空间的最小生成元长度与秩的定义直接推出。证毕。
\end{proof}

\begin{proposition}[耦合参数的主因子定律]\label{prop:xi-symq-rankone-coupling-t-power}
\leanverified{paper\_xi\_symq\_rankone\_coupling\_t\_power}
对任意 \(t\in\CC\)，定义
\[
p_{q,t}(z):=\det\!\bigl(zI_{q+1}-(B_q+tE_q)\bigr).
\]
则对一切 \(q\ge 2\) 有
\[
\operatorname{Res}_z(g_q,p_{q,t})
=t^{\,q+1}\operatorname{Res}_z(g_q,p_{q,1}).
\]
\end{proposition}
\begin{proof}
沿用定理 \ref{thm:xi-symq-kj-resultant-fibonacci} 的特征基与 \(u\) 的记号，有
\(
B_q+tE_q=\mathrm{diag}(\mu_i)+tuu^{\mathsf T}
\)。
故对每个 \(\mu_i\) 成立 \(p_{q,t}(\mu_i)=-t\,u_i^2g_q'(\mu_i)\)，从而
\[
\operatorname{Res}_z(g_q,p_{q,t})
=\prod_{i=0}^{q}p_{q,t}(\mu_i)
=t^{q+1}\prod_{i=0}^{q}p_{q,1}(\mu_i)
=t^{q+1}\operatorname{Res}_z(g_q,p_{q,1}).
\]
证毕。
\end{proof}

\begin{proposition}[谱差积的 \(p\)-进估值公式：数字和项与 Fibonacci 赋值项]\label{prop:xi-symq-resultant-padic-valuation}
\leanverified{paper\_xi\_symq\_resultant\_padic\_valuation}
令
\[
R_q:=\bigl|\operatorname{Res}_z(g_q,p_q)\bigr|
=\left(\prod_{i=0}^{q}\binom{q}{i}\right)\left(\prod_{k=1}^{q}F_k^{\,2(q+1-k)}\right).
\]
则对任意素数 \(p\) 有
\[
\nu_p(R_q)=\sum_{i=0}^{q}\nu_p\!\binom{q}{i}
\;+\;\sum_{k=1}^{q}2(q+1-k)\,\nu_p(F_k),
\]
并且
\[
\sum_{i=0}^{q}\nu_p\!\binom{q}{i}
=\frac{2\sum_{n=0}^{q}s_p(n)-(q+1)s_p(q)}{p-1},
\]
其中 \(s_p(n)\) 表示 \(n\) 的 \(p\)-进展开的数字和。
\end{proposition}
\begin{proof}
第一式由 \(R_q\) 的乘积结构与 \(p\)-进估值的可加性直接得到。
第二式由 Legendre 公式
\(
\nu_p(n!)=\frac{n-s_p(n)}{p-1}
\)
推出
\[
\nu_p\binom{q}{i}
=\nu_p(q!)-\nu_p(i!)-\nu_p((q-i)!)
=\frac{s_p(i)+s_p(q-i)-s_p(q)}{p-1},
\]
对 \(i=0,\dots,q\) 求和并利用 \(\sum_{i=0}^{q}s_p(q-i)=\sum_{i=0}^{q}s_p(i)\) 即得。证毕。
\end{proof}

\begin{proposition}[判别式/结果式比值的纯 \(5\)-幂缺陷]\label{prop:xi-symq-disc-over-resultant-pure-5-defect}
对一切 \(q\ge 2\) 成立
\[
\frac{\Disc(g_q)}{\operatorname{Res}_z(g_q,p_q)}
=(-1)^{\frac{(q+1)(q+2)}2}\cdot
\frac{5^{\frac{q(q+1)}2}}{\displaystyle\prod_{i=0}^{q}\binom{q}{i}}.
\]
\end{proposition}
\begin{proof}
由定理 \ref{thm:xi-symq-k-discriminant-fibonacci-vandermonde} 与定理 \ref{thm:xi-symq-kj-resultant-fibonacci} 直接相除即可；
其中 \(\prod_{k=1}^{q}F_k^{\,2(q+1-k)}\) 在该比值中严格消去。证毕。
\end{proof}

\endinput


\begin{theorem}[Perron 根的纯 Fibonacci 固定点方程]\label{thm:xi-replica-softcore-perron-fibonacci-fixed-point}
记 \(\rho_q:=\rho(A^{(q)})\) 为 \(\left.T^{(q)}\right|_{\mathcal H_q}\) 的 Perron 特征值。
则存在唯一实数 \(x_q>1\) 使
\[
\rho_q=2^{q-1}x_q,
\]
并且 \(x_q\) 满足严格收敛的固定点方程
\begin{equation}\label{eq:xi-replica-softcore-perron-fixed-point}
x_q
=1+\sum_{m=1}^{\infty}\left(\frac{F_{m+3}}{2^{m+1}}\right)^{q}\,x_q^{-m}.
\end{equation}
该刻画在每个 \(q\ge 2\) 下给出唯一正根。
\end{theorem}
\begin{proof}
由定理 \ref{thm:xi-replica-softcore-exceptional-weighted-secular}，Perron 根 \(\rho_q\) 为割线方程
\begin{equation}\label{eq:xi-replica-softcore-perron-secular}
\sum_{i=0}^{q}\frac{w_i^{(q)}}{\rho_q-d_i}=2
\end{equation}
在顶端区间的唯一解，其中 \(d_i=\mu_i^{(q)}/2\)。
并且由 Rayleigh 商下界（取向量 \(u=D\mathbf 1\) 于相似对称表示 \(\widetilde A^{(q)}\) 上计算）可得
\(
\rho_q\ge (4^q+3^q)/2^{q+1}
=2^{q-1}+3^q/2^{q+1}
\)；
从而 \(\rho_q>\max_{0\le i\le q}d_i=\varphi^q/2\)，并对每个 \(i\) 有绝对收敛展开
\[
\frac{1}{\rho_q-d_i}
=\frac{1}{\rho_q}\sum_{m=0}^{\infty}\left(\frac{d_i}{\rho_q}\right)^{m}.
\]
将其代入 \eqref{eq:xi-replica-softcore-perron-secular} 并由一致收敛交换求和次序，得到
\[
2
=\frac{1}{\rho_q}\sum_{m=0}^{\infty}\rho_q^{-m}\sum_{i=0}^{q} w_i^{(q)}\,d_i^{\,m}.
\]
由定理 \ref{thm:xi-replica-softcore-fibonacci-moment-collapse}，
\(
\sum_i w_i^{(q)}d_i^{\,m}=2^{-m}F_{m+3}^q
\)，故
\[
2
=\frac{1}{\rho_q}\sum_{m=0}^{\infty}\rho_q^{-m}\,2^{-m}F_{m+3}^{\,q}.
\]
令 \(\rho_q=2^{q-1}x_q\)，并将 \(m=0\) 项（\(F_3=2\)）分离，即得
\[
2
=\frac{2}{x_q}+ \frac{2}{x_q}\sum_{m=1}^{\infty}\left(\frac{F_{m+3}}{2^{m+1}}\right)^{q}x_q^{-m}.
\]
两边同乘 \(x_q/2\) 得到固定点方程 \eqref{eq:xi-replica-softcore-perron-fixed-point}。

右端作为 \(x\) 的函数在 \((1,\infty)\) 上严格递减，而左端 \(x\mapsto x\) 严格递增，故解唯一。证毕。
\end{proof}

\begin{corollary}[Perron 根的两项主导展开与可控余项]\label{cor:xi-replica-softcore-perron-two-term-expansion}
令
\[
b_m:=\frac{F_{m+3}}{2^{m+1}}\qquad(m\ge 1),
\qquad
\varepsilon_q:=\sum_{m=1}^{\infty}b_m^{\,q}.
\]
则 \(x_q\in(1,1+\varepsilon_q)\)，并且
\begin{equation}\label{eq:xi-replica-softcore-perron-two-term-expansion}
\rho_q
=2^{q-1}+\frac{3^q}{2^{q+1}}+R_q,
\end{equation}
其中余项满足显式上界
\begin{equation}\label{eq:xi-replica-softcore-perron-remainder-bound}
0<R_q
<2^{q-1}\sum_{m=2}^{\infty}b_m^{\,q}
\le 2^{q-1}\cdot \frac{b_2^{\,q}}{1-(\varphi/2)^q}
=\frac{5^q}{2^{2q+1}}\cdot\frac{1}{1-(\varphi/2)^q}.
\end{equation}
因此
\(
\rho_q
=2^{q-1}+\frac{3^q}{2^{q+1}}+O\!\bigl(5^q/2^{2q}\bigr)
\)
当 \(q\to\infty\)。
\end{corollary}
\begin{proof}
由定理 \ref{thm:xi-replica-softcore-perron-fibonacci-fixed-point}，
\[
x_q-1=\sum_{m=1}^{\infty} b_m^{\,q}x_q^{-m}.
\]
因 \(x_q>1\)，即得 \(0<x_q-1<\sum_{m\ge 1}b_m^q=\varepsilon_q\)，从而 \(x_q\in(1,1+\varepsilon_q)\)。
将首项 \(m=1\) 分离并乘以 \(2^{q-1}\) 得
\[
\rho_q
=2^{q-1}+2^{q-1}\frac{b_1^{\,q}}{x_q}
2^{q-1}\sum_{m=2}^{\infty}b_m^{\,q}x_q^{-m}.
\]
由于 \(b_1=F_4/2^2=3/4\) 且 \(x_q>1\)，有
\(
2^{q-1}b_1^q/x_q<2^{q-1}b_1^q=3^q/2^{q+1}
\)。
将差额并入余项并注意最后一项为正，得到展开式 \eqref{eq:xi-replica-softcore-perron-two-term-expansion}，且
\[
0<R_q<2^{q-1}\sum_{m=2}^{\infty}b_m^{\,q}x_q^{-m}\le 2^{q-1}\sum_{m=2}^{\infty}b_m^{\,q}.
\]
再由 Fibonacci 比值单调性 \(F_{n+1}/F_n<\varphi\)（\(n\ge 1\)）推出
\[
\frac{b_{m+1}}{b_m}=\frac{F_{m+4}}{2F_{m+3}}<\frac{\varphi}{2},
\]
故 \(\sum_{m\ge 2}b_m^q\) 由等比级数控制，得到 \eqref{eq:xi-replica-softcore-perron-remainder-bound}。证毕。
\end{proof}

\begin{theorem}[Fibonacci 幂 Hankel 行列式的完全乘积公式]\label{thm:xi-replica-softcore-fibonacci-power-hankel-determinant}
定义 \((q+1)\times(q+1)\) Hankel 矩阵
\[
H_q:=\bigl(F_{i+j+3}^{\,q}\bigr)_{0\le i,j\le q}.
\]
则对一切 \(q\ge 1\) 成立精确闭式
\begin{equation}\label{eq:xi-replica-softcore-fibonacci-power-hankel-determinant}
\det(H_q)
=\left(\prod_{i=0}^{q}\binom{q}{i}\right)
\left(\prod_{k=1}^{q}F_k^{\,2(q+1-k)}\right).
\end{equation}
右端为整数的纯乘积表达，且不含任何代数数域延拓。
\end{theorem}
\begin{proof}
由 Binet 形式与 \(\alpha=\varphi^3/\sqrt5\)、\(\beta=-\psi^3/\sqrt5\)，有
\(
F_{n+3}=\alpha\varphi^n+\beta\psi^n
\)；
从而对任意 \(n\ge 0\),
\[
F_{n+3}^{\,q}
=\sum_{i=0}^{q}\binom{q}{i}\alpha^{q-i}\beta^{i}(\varphi^{q-i}\psi^{i})^{n}
=\sum_{i=0}^{q}c_i r_i^{\,n},
\]
其中 \(r_i:=\varphi^{q-i}\psi^{i}\) 且 \(c_i:=\binom{q}{i}\alpha^{q-i}\beta^{i}\)。
令 \(V:=(r_i^{\,j})_{0\le i,j\le q}\) 为 Vandermonde 型矩阵，并令 \(C:=\mathrm{diag}(c_0,\dots,c_q)\)。
则 \(H_q=V^{\mathsf T}CV\)，从而
\begin{equation}\label{eq:xi-replica-softcore-hankel-det-vandermonde-step}
\det(H_q)=(\det V)^2\prod_{i=0}^{q}c_i
=\left(\prod_{0\le i<j\le q}(r_j-r_i)^2\right)\left(\prod_{i=0}^{q}c_i\right).
\end{equation}

先算系数乘积。
由 \(\sum_{i=0}^{q} i=\sum_{i=0}^{q}(q-i)=q(q+1)/2\) 与
\(
\alpha\beta=(\varphi^3/\sqrt5)(-\psi^3/\sqrt5)=-(\varphi\psi)^3/5=1/5
\),
得到
\[
\prod_{i=0}^{q}c_i
=\left(\prod_{i=0}^{q}\binom{q}{i}\right)(\alpha\beta)^{q(q+1)/2}
=\left(\prod_{i=0}^{q}\binom{q}{i}\right)5^{-q(q+1)/2}.
\]

再算 Vandermonde 因子。
注意 \(r_i=\varphi^q(\psi/\varphi)^i=\varphi^q r^i\)，其中 \(r:=\psi/\varphi=-\varphi^{-2}\)。
于是
\[
\prod_{0\le i<j\le q}(r_j-r_i)^2
=(\varphi^q)^{q(q+1)}\prod_{0\le i<j\le q}(r^j-r^i)^2.
\]
而几何节点的 Vandermonde 恒等式给出
\[
\prod_{0\le i<j\le q}(r^j-r^i)^2
=\prod_{k=1}^{q}(r^k-1)^{2(q+1-k)}\cdot r^{\,2\sum_{0\le i<j\le q}i}.
\]
对每个 \(k\ge 1\)，有
\[
1-r^k
=1-(-\varphi^{-2})^k
=\frac{\varphi^{2k}-(-1)^k}{\varphi^{2k}}
=\frac{\varphi^k(\varphi^k-\psi^k)}{\varphi^{2k}}
=\frac{\sqrt5\,F_k}{\varphi^k},
\]
故
\(
(r^k-1)^2=(1-r^k)^2=5F_k^2/\varphi^{2k}
\)。
将其代回并整理 \(\varphi\) 的幂次，可验证所有 \(\varphi\) 因子完全相消，从而得到
\[
\prod_{0\le i<j\le q}(r_j-r_i)^2
=5^{q(q+1)/2}\prod_{k=1}^{q}F_k^{\,2(q+1-k)}.
\]
与 \eqref{eq:xi-replica-softcore-hankel-det-vandermonde-step} 中的 \(5^{-q(q+1)/2}\) 系数因子精确抵消，得到 \eqref{eq:xi-replica-softcore-fibonacci-power-hankel-determinant}。证毕。
\end{proof}

\begin{corollary}[交错几何点的 Vandermonde 平方塌缩为 Fibonacci 乘积]\label{cor:xi-replica-softcore-alternating-geom-vandermonde}
令
\[
\xi_i:=(-1)^i\varphi^{\,q-2i}\qquad(0\le i\le q).
\]
则成立严格恒等式
\begin{equation}\label{eq:xi-replica-softcore-alternating-geom-vandermonde}
\prod_{0\le i<j\le q}(\xi_j-\xi_i)^2
=5^{q(q+1)/2}\prod_{k=1}^{q}F_k^{\,2(q+1-k)}\in\ZZ.
\end{equation}
因此左端虽表面上属于 \(\QQ(\sqrt5)\) 的代数整数乘积，其平方判别式却在 \(\ZZ\) 中塌缩为完全显式的 Fibonacci 乘积。
\end{corollary}
\begin{proof}
在定理 \ref{thm:xi-replica-softcore-fibonacci-power-hankel-determinant} 的证明中，节点 \(r_i=\varphi^{q-i}\psi^i\) 满足
\(
r_i=(-1)^i\varphi^{q-2i}=\xi_i
\)。
因此其中已得的 Vandermonde 平方闭式即为 \eqref{eq:xi-replica-softcore-alternating-geom-vandermonde}。证毕。
\end{proof}

\begin{theorem}[严格交错与全矩阵惯性平衡]\label{thm:xi-replica-softcore-full-inertia-balance}
设
\[
d_i:=\frac{\mu_i^{(q)}}2\qquad(0\le i\le q),
\]
按降序重排为
\(
d_{(0)}>d_{(1)}>\cdots>d_{(q)}
\)。
则例外谱可重排为
\(
\lambda_{(0)}>\lambda_{(1)}>\cdots>\lambda_{(q)}
\)
并满足严格交错
\[
\lambda_{(0)}>d_{(0)}>\lambda_{(1)}>d_{(1)}>\cdots>\lambda_{(q)}>d_{(q)}.
\]
进一步有：
\begin{enumerate}
  \item 例外谱负特征值个数为 \(\lfloor(q+1)/2\rfloor\)，正特征值个数为 \(\lceil(q+1)/2\rceil\)；
  \item 全矩阵惯性满足
  \[
  \operatorname{In}\!\bigl(T^{(q)}\bigr)=\bigl(2^{q-1},\,2^{q-1},\,0\bigr).
  \]
\end{enumerate}
\end{theorem}
\begin{proof}
交错性由定理 \ref{thm:xi-replica-softcore-exceptional-weighted-secular} 的割线函数直接推出：
在每个相邻极点区间与顶端区间，方程恰有一根。
例外谱符号计数与零点排除亦可由推论 \ref{cor:xi-replica-softcore-exceptional-inertia} 直接得到。

体谱中奇 \(i\) 给出负特征值，偶 \(i\) 给出正特征值，故
\[
N_-^{\mathrm{bulk}}
=
\sum_{\substack{0\le i\le q\\ i\ \mathrm{odd}}}
\left(\binom qi-1\right)
=
2^{q-1}-\left\lfloor\frac{q+1}{2}\right\rfloor,
\]
\[
N_+^{\mathrm{bulk}}
=
\sum_{\substack{0\le i\le q\\ i\ \mathrm{even}}}
\left(\binom qi-1\right)
=
2^{q-1}-\left\lceil\frac{q+1}{2}\right\rceil.
\]
再加上例外谱计数
\[
N_-^{\mathrm{exc}}=\left\lfloor\frac{q+1}{2}\right\rfloor,\qquad
N_+^{\mathrm{exc}}=\left\lceil\frac{q+1}{2}\right\rceil,
\]
得到
\[
N_-=2^{q-1},\qquad N_+=2^{q-1}.
\]
由前述零点排除，\(N_0=0\)。故惯性公式成立。证毕。
\end{proof}

\begin{proposition}[体谱特征空间的零和超平面方程]\label{prop:xi-replica-softcore-bulk-eigenspace-zero-sum-hyperplane}
取 \(K\) 的单位正交本征向量 \(e_+,e_-\)，并对任意 \(S\subseteq[q]\) 定义
\[
v_S:=\bigotimes_{r=1}^{q}\xi_r,\qquad
\xi_r=
\begin{cases}
e_-,&r\in S,\\
e_+,&r\notin S.
\end{cases}
\]
当 \(|S|=i\) 时，\(\{v_S\}\) 构成 \(E_i\) 的一组正交基。
并且
\[
\ker\!\left(T^{(q)}-\frac{\mu_i^{(q)}}2I\right)
=
\left\{
\sum_{|S|=i}a_Sv_S:\ 
\sum_{|S|=i}a_S\langle v_S,\mathbf1\rangle=0
\right\}
=E_i\cap\mathbf1^\perp,
\]
其维数为 \(\binom qi-1\)。
特别地，由 \(\langle v_S,\mathbf1\rangle\) 仅依赖于 \(|S|=i\) 且非零，上式等价为
\[
\sum_{|S|=i}a_S=0.
\]
\end{proposition}
\begin{proof}
张量积谱给出 \(\{v_S:|S|=i\}\) 为 \(E_i\) 的正交基。
对任意 \(x\in E_i\)，
\[
T^{(q)}x=\frac12\bigl(D^{(q)}x+J_qx\bigr)=\frac{\mu_i^{(q)}}2x+\frac12J_qx.
\]
因此
\[
\left(T^{(q)}-\frac{\mu_i^{(q)}}2I\right)x=0
\iff J_qx=0
\iff x\perp\mathbf1.
\]
写 \(x=\sum_{|S|=i}a_Sv_S\) 即得线性约束式。
由于 \(\mathbf1\) 在 \(E_i\) 上投影非零，该约束为一条独立线性方程，故维数为 \(\binom qi-1\)。证毕。
\end{proof}

\endinput
